\section{Derive the resolvent of a second-order Jordan block}
\label{app:jordan-resolvent}
% BEGIN CURATED SUBJECT INDEX
\index{resolvent}
\index{exceptional point}
\index{Jordan chain}
% END CURATED SUBJECT INDEX

On the two-dimensional root subspace, write
\begin{equation}
  J=\vsub{E}{EP}\identity+N,
  \qquad N^2=0,
  \qquad N\ne0 .
  \label{eq:jordan-nilpotent}
\end{equation}
Then
\begin{equation}
  (z-J)^{-1}
  =\frac{\identity}{z-\vsub{E}{EP}}
  +\frac{N}{(z-\vsub{E}{EP})^2} .
  \label{eq:jordan-resolvent}
\end{equation}
The double-pole coefficient is nilpotent.  Under a generic perturbation,
$N$ and the perturbation matrix combine to split the double pole into the two
square-root branches in Eq.~\eqref{eq:ep-puiseux}.  Individual rank-one
projectors diverge as the EP is approached, but their sum tends to the finite
Riesz projector onto the root subspace.  Numerically, it is therefore more
stable to track the two-dimensional contour-integral projector than two
separately normalized eigenvectors.
